\documentclass[../Main.tex]{subfiles}
\begin{document}
  \chapter{2018-2019学年微积分（一）（上）期末考试}

  \section{单项选择题(每小题 3 分，共 18 分)}
  \begin{enumerate}
    \item 以下关于数列的命题，正确的是【\quad】.

      \begin{tabular}{l}
        \textbf{A.} 一个有界数列与一个无界数列的和是无界数列  \\
        \textbf{B.} 两个无界数列的和是无界数列         \\
        \textbf{C.} 一个有界数列与一个无界数列的乘积是无界数列 \\
        \textbf{D.} 两个无界数列的乘积是无界数列
      \end{tabular}
      \vspace{1em}

    \item 设函数$f(x)$在$x=a$处可导，则函数$\sqrt[3]{f(x)}$在$x=a$处【\quad】.

      \begin{tabular}{l@{\hspace{2cm}}l@{\hspace{2cm}}l@{\hspace{2cm}}l}
        \textbf{A.} 可导 & \textbf{B.} 不连续 & \textbf{C.} 连续但不一定可导 & \textbf{D.} 不可导
      \end{tabular}
      \vspace{1em}

    \item 设函数$f(x)$在区间$(a,b)$内连续，则$f(x)$在区间$(a,b)$内【\quad】.

      \begin{tabular}{l@{\hspace{2cm}}l@{\hspace{2cm}}l@{\hspace{2cm}}l}
        \textbf{A.} 有界 & \textbf{B.} 可导 & \textbf{C.} 存在最大值 & \textbf{D.} 原函数存在
      \end{tabular}
      \vspace{1em}

    \item 函数$f(x)=x^{4}-2x^{3}$有【\quad】.

      \begin{tabular}{l@{\hspace{2cm}}l}
        \textbf{A.} 一个极小值和一个极大值 & \textbf{B.} 一个极小值 \\
        \textbf{C.} 两个极小值       & \textbf{D.} 两个极大值
      \end{tabular}
      \vspace{1em}

    \item 设函数$f(x)$在区间$(a,b)$内满足$f'(x)<0$，$f''(x)>0$，则在区间$(a,b)$内【\quad】.

      \begin{tabular}{l@{\hspace{2cm}}l}
        \textbf{A.} $f(x)$单调减少，曲线$y=f(x)$下凸 & \textbf{B.} $f(x)$单调减少，曲线$y=f(x)$上凸 \\
        \textbf{C.} $f(x)$单调增加，曲线$y=f(x)$下凸 & \textbf{D.} $f(x)$单调增加，曲线$y=f(x)$上凸
      \end{tabular}
      \vspace{1em}

    \item 设$M=\int_{0}^{\frac{\pi}{4}}\sqrt{\sin x}\mathrm{d}x$，$N=\int_{0}^{\frac{\pi}{4}}
      \sqrt{\sec x}\mathrm{d}x$，$K=\int_{0}^{\frac{\pi}{4}}\sqrt{\tan x}\mathrm{d}
      x$，则$M,N,K$的大小关系为【\quad】.

      \begin{tabular}{l@{\hspace{1.5cm}}l@{\hspace{1.5cm}}l@{\hspace{1.5cm}}l}
        \textbf{A.} $M<N<K$ & \textbf{B.} $M<K<N$ & \textbf{C.} $N<M<K$ & \textbf{D.} $K<N<M$
      \end{tabular}
      \vspace{1em}
  \end{enumerate}

  \section{填空题(每小题 4 分，共 16 分)}
  \begin{enumerate}
    \item[7.] 设$u=x+a\ln(1-x)+bx\sin(3x)$是$x$的$3$阶无穷小，则$a=\underline{\hspace{2cm}
      }$，$b=\underline{\hspace{2cm}}$.
      \vspace{0.5em}

    \item[8.] 曲线$y=x^{3}-6x^{2}+5x+3$的拐点坐标为$\underline{\hspace{3cm}}$.
      \vspace{0.5em}

    \item[9.] $\lim_{x\to0}\frac{\displaystyle\int_{0}^{1-\cos x}\tan t\mathrm{d}t}{\sin
      x^{4}}=\underline{\hspace{3cm}}$.
      \vspace{0.5em}

    \item[10.] 曲线$x=2\cos^{3} t$，$y=2\sin^{3} t$，$0\leq t \leq 2\pi$的长度为$\underline
      {\hspace{3cm}}$.
  \end{enumerate}

  \section{计算题(每小题 7 分，共 42 分)}
  \begin{enumerate}
    \item[11.] 求曲线$y=\sqrt{4x^{2}+2x+5}(x>0)$的渐近线.
      \vspace{12em}

    \item[12.] 写出$f(x)=\ln(1+x)$带Lagrange余项的$n$阶Maclaurin公式.
      \vspace{12em}

    \item[13.] 求不定积分$I=\int \frac{\sqrt{x-3}}{2x}\mathrm{d}x$.
      \vspace{12em}

    \item[14.] 求定积分$I=\int_{0}^{\frac{1}{2}}x\arcsin x\mathrm{d}x$.
      \vspace{11em}

    \item[15.] 求反常积分$I= \int_{0}^{+\infty}\mathrm{e}^{-2x}\sin x\mathrm{d}x$.
      \vspace{12em}

    \item[16.] 求微分方程$y'+\frac{y}{x}=y^{2}\ln x$的通解.
      \vspace{12em}
  \end{enumerate}

  \section{综合题(每小题 7 分，共 14 分)}
  \begin{enumerate}
    \item[17.] 设函数$f(x)=
      \begin{cases}
        (1+x)^{\frac{1}{x}}, & x\in(-1,0)\cup(0,+\infty), \\
        \mathrm{e},          & x=0
      \end{cases}$，求$f'(x)$并讨论$f'(x)$的连续性.
      \vspace{12em}

    \item[18.] 求平面图形$0\leq y\leq \cos x$，$0\leq x\leq \frac{\pi}{2}$绕$y$轴旋转所得立体的体积.
      \vspace{15em}
  \end{enumerate}

  \section{证明题(每小题 5 分，共 10 分)}
  \begin{enumerate}
    \item[19.] 设函数$f(x)$在区间$[a,b]$上二阶可导，且$f(a)=f(b)$，$|f''(x)|\leq M$，证明：
      \[
        |f'(a)+f'(b)|\leq M(b-a).
      \]
      \vspace{15em}

    \item[20.] 设$f(x)$在区间$[a,b]$上二阶可导，并且$f''(x)\leq 0$，证明：
      \[
        \int_{a}^{b}f(x)\mathrm{d}x \geq \frac{(b-a)(f(a)+f(b))}{2}.
      \]
      \vspace{15em}
  \end{enumerate}

  \chapter{2018-2019学年微积分（一）（上）期末考试参考答案}
  \section{单项选择题 (每小题 3 分，共 18 分)}
  \begin{enumerate}
    \item \textbf{Solution}. A.

      对于A选项，假设$\{a_{n}\}$是一个有界数列，$\{b_{n}\}$是一个无界数列，

      即$\forall\,n\in\mathbf{N}^{+}$，$\exists\,M>0$，使得$|a_{n}|\leq M$，又$\forall
      \,N>0$，$\exists\,n_{0}\in\mathbf{N}^{+}$，使得$|b_{n_0}|>N$.

      则$|a_{n_0}+b_{n_0}|\geq |b_{n_0}|-|a_{n_0}|\geq N-M$，$\forall\,N-M>0$，所以$\{
      a_{n}+b_{n}\}$是一个无界数列.

      对于B选项，考虑$a_{n}= n$，$b_{n}=-n$，则$a_{n}+b_{n}\equiv 0$，$\{a_{n}+b_{n}
      \}$是一个有界数列.

      对于C选项，考虑$a_{n}=0$，$b_{n}=n$，则$a_{n}b_{n}\equiv 0$，$\{a_{n}b_{n}\}$是一个有界数列.

      对于D选项，考虑$a_{n}=
      \begin{cases}
        n, & n=2k,  \\
        0, & n=2k-1
      \end{cases}$，$b_{n}=
      \begin{cases}
        0, & n=2k,  \\
        n, & n=2k-1
      \end{cases}$，其中$k\in\mathbf{N}^{+}$，

      则$a_{n}b_{n}\equiv 0$，$\{a_{n}b_{n}\}$是一个有界数列.

    \item \textbf{Solution}. C.

      由题意可知$\lim_{x\to a}\frac{f(x)-f(a)}{x-a}$存在，所以$\lim_{x\to a}\left
      (f(x)-f(a)\right)=\lim_{x\to a}\frac{f(x)-f(a)}{x-a}\cdot \lim_{x\to a}(x-a
      )=0$，

      即$\lim_{x\to a}f(x)=f(a)$，所以$\lim_{x\to a}\sqrt[3]{f(x)}=\sqrt[3]{\lim_{x\to
      a}f(x)}=\sqrt[3]{f(a)}$，$\sqrt[3]{f(x)}$在$x=a$处连续.

      考虑$f(x)=x$，$a=0$，则$\sqrt[3]{f(x)}=\sqrt[3]{x}$在$x=0$处不可导. 所以$\sqrt[3]{f(x)}$在$x
      =a$处连续但不一定可导.

    \item \textbf{Solution}. D.

      对于A选项，考虑$f(x)=\frac{1}{x}$，取$a=-1$，$b=1$，则$f(x)$在区间$(a,b)$内连续但无界.

      对于B，C选项，考虑$f(x)=|x|$，取$a=-1$，$b=1$，则$f(x)$在区间$(a,b)$内连续但不存在最大值，

      并且存在不可导点$x=0$.

      连续函数在区间内必然可积，所以D选项正确.

    \item \textbf{Solution}. B.

      $f'(x)=4x^{3}-6x^{2}=2x^{2}(2x-3)$，

      所以函数$f(x)$在区间$\left(-\infty,\frac{3}{2}\right)$上单调递减，在区间$\left
      (\frac{3}{2},+\infty\right)$上单调递增.

      函数$f(x)$只有一个唯一的极小值点$x=\frac{3}{2}$，没有极大值点.

    \item \textbf{Solution}. A.

      由$f'(x)<0$，$f''(x)>0$可知函数$f(x)$在区间$(a,b)$内单调递减，曲线$y=f(x)$下凸.

    \item \textbf{Solution}. B.

      在区间$\left[0,\frac{\pi}{4}\right]$内，$\sin x\leq \tan x\leq \sec x$，所以$M<K<N$.
  \end{enumerate}

  \section{填空题(每小题 4 分，共 16 分)}
  \begin{enumerate}
    \item[7.] \textbf{Solution}. $a=1$，$b=\frac{1}{6}$.

      由Taylor公式，
      \[
        \begin{aligned}
          u = x+a\ln(1-x)+bx\sin(3x) = & x-a\left(x+\frac{x^2}{2}+\frac{x^3}{3}+o(x^{3})\right)+bx\left(3x+o(x^{2})\right) \\
          =                            & (1-a)x+\left(-\frac{a}{2}+3b\right)x^{2}-\frac{a}{3}x^{3}+o(x^{3}).
        \end{aligned}
      \]

      所以$1-a=0$，$-\frac{a}{2}+3b=0$，解得$a=1$，$b=\frac{1}{6}$.

    \item[8.] \textbf{Solution}. $(2,-3)$.

      $y'=3x^{2}-12x+5$，$y''=6x-12$，令$y''=0$得$x=2$，且当$x<2$时，$y''<0$，当$x
      >2$时，$y''>0$，

      所以曲线$y=x^{3}-6x^{2}+5x+3$具有拐点$(2,-3)$.

    \item[9.] \textbf{Solution}. $\frac{1}{8}$.
      \[
        \begin{aligned}
          \lim_{x\to0}\frac{\displaystyle\int_{0}^{1-\cos x}\tan t\mathrm{d}t}{\sin x^4}= & \lim_{x\to0}\frac{\tan(1-\cos x)\cdot \sin x}{4x^3} \\
          =                                                                               & \lim_{x\to0}\frac{1-\cos x}{4x^2}                   \\
          =                                                                               & \lim_{x\to0}\frac{\sin x}{8x}= \frac{1}{8}.
        \end{aligned}
      \]

    \item[10.] \textbf{Solution}. $12$.

      $\mathrm{d}s=\sqrt{x'^{2}(t)+y'^{2}(t)}\mathrm{d}t=\sqrt{(-6\cos^{2} t\sin
      t)^{2}+(6\sin^{2} t\cos t)^{2}}\mathrm{d}t=6\sqrt{\sin^{2} t\cos^{2} t}\mathrm{d}
      t=3|\sin 2t|\mathrm{d}t$.

      故$s=3\int_{0}^{2\pi}|\sin 2t|\mathrm{d}t=6\int_{0}^{\pi}|\sin u|\mathrm{d}
      u=12\int_{0}^{\frac{\pi}{2}}\sin u\mathrm{d}u=12$.
  \end{enumerate}
  \section{计算题(每小题 7 分，共 42 分)}
  \begin{enumerate}
    \item[11.] \textbf{Solution}. 曲线没有竖直渐近线.

      $\lim_{x\to+\infty}\frac{y}{x}=\lim_{x\to+\infty}\frac{\sqrt{4x^{2}+2x+5}}{x}
      =\lim_{x\to+\infty}\sqrt{4+\frac{2}{x}+\frac{5}{x^{2}}}=2$，

      $\begin{aligned}
        \lim_{x\to+\infty}(y-2x)= & \lim_{x\to+\infty}\left(\sqrt{4x^2+2x+5}-2x\right)                                        \\
        =                         & \lim_{x\to+\infty}\frac{2x+5}{\sqrt{4x^2+2x+5}+2x}                                        \\
        =                         & \lim_{x\to+\infty}\frac{2+\frac{5}{x}}{\sqrt{4+\frac{2}{x}+\frac{5}{x^2}}+2}=\frac{1}{2}.
      \end{aligned}$

      所以曲线的斜渐近线为$y=2x+\frac{1}{2}$.

    \item[12.] \textbf{Solution}. $f(x)=x-\frac{x^{2}}{2}+\frac{x^{3}}{3}-\cdots +
      \frac{(-1)^{n-1}}{n}x^{n}+R_{n}(x)$，

      其中$R_{n}(x)=\frac{f^{(n+1)}(\xi)}{(n+1)!}x^{n+1}$，$\xi$介于$0$与$x$之间.

      由于$f^{(n+1)}(x)=\frac{(-1)^{n}\,n!}{(1+x)^{n+1}}$，所以$R_{n}(x)=\frac{(-1)^{n}\,x^{n+1}}{(n+1)(1+\xi)^{n+1}}$.

    \item[13.] \textbf{Solution}. 令$u=\sqrt{x-3}$，则$x=t^{2}+3$，$\mathrm{d}x=2
      u\mathrm{d}u$，
      \[
        \begin{aligned}
          I = & \int \frac{\sqrt{x-3}}{2x}\mathrm{d}x = \int \frac{u}{2(u^2+3)}\cdot 2u\mathrm{d}u                         \\
          =   & \int \left(1-\frac{3}{u^2+3}\right)\mathrm{d}u                                                             \\
          =   & u-3\cdot\frac{1}{\sqrt{3}}\arctan\frac{u}{\sqrt{3}}+ C = \sqrt{x-3}-\sqrt{3}\arctan\sqrt{\frac{x-3}{3}}+C.
        \end{aligned}
      \]

    \item[14.] \textbf{Solution}.
      \[
        \begin{aligned}
          I = \frac{1}{2}\int_{0}^{\frac{1}{2}}\arcsin x\mathrm{d}x^{2} = & \frac{1}{2}x^{2}\arcsin x\bigg|_{0}^{\frac{1}{2}}-\frac{1}{2}\int_{0}^{\frac{1}{2}}\frac{x^2}{\sqrt{1-x^2}}\mathrm{d}x  \\
          =                                                               & \frac{1}{8}\cdot\frac{\pi}{6}-\frac{1}{2}\int_{0}^{\frac{\pi}{6}}\frac{\sin^2 t}{\cos t}\cdot \cos t\mathrm{d}t         \\
          =                                                               & \frac{\pi}{48}-\frac{1}{4}\int_{0}^{\frac{\pi}{6}}(1-\cos 2t)\mathrm{d}t                                                \\
          =                                                               & \frac{\pi}{48}-\frac{1}{4}\left(\frac{\pi}{6}-\frac{1}{2}\sin\frac{\pi}{3}\right) = \frac{\sqrt{3}}{16}-\frac{\pi}{48}.
        \end{aligned}
      \]

    \item[15.] \textbf{Solution}. 记$J = \int \mathrm{e}^{-2x}\sin x\mathrm{d}x$，则
      \[
        \begin{aligned}
          J = -\int \mathrm{e}^{-2x}\mathrm{d}\cos x = & -\mathrm{e}^{-2x}\cos x - 2\int \mathrm{e}^{-2x}\cos x\mathrm{d}x                           \\
          =                                            & -\mathrm{e}^{-2x}\cos x - 2\int \mathrm{e}^{-2x}\mathrm{d}\sin x                            \\
          =                                            & -\mathrm{e}^{-2x}\cos x - 2\mathrm{e}^{-2x}\sin x - 4\int \mathrm{e}^{-2x}\sin x\mathrm{d}x \\
          =                                            & -\mathrm{e}^{-2x}\cos x - 2\mathrm{e}^{-2x}\sin x - 4J.
        \end{aligned}
      \]

      所以$J = -\frac{1}{5}\mathrm{e}^{-2x}\cos x - \frac{2}{5}\mathrm{e}^{-2x}\sin
      x + C$，

      $I=\int_{0}^{+\infty}\mathrm{e}^{-2x}\sin x\mathrm{d}x=\left[-\frac{1}{5}\mathrm{e}
      ^{-2x}\cos x - \frac{2}{5}\mathrm{e}^{-2x}\sin x\right]_{0}^{+\infty}=\frac{1}{5}$.

    \item[16.] \textbf{Solution}. 令$u=\frac{1}{y}$，则$u'=-\frac{y'}{y^{2}}$，所以$y
      '=-\frac{u'}{u^{2}}$.

      微分方程$y'+\frac{y}{x}=y^{2}\ln x$变形为$-\frac{u'}{u^{2}}+\frac{1}{xu}=\frac{\ln
      x}{u^{2}}$，即$u'-\frac{1}{x}u=-\ln x$.

      由一阶非齐次线性微分方程的通解公式得
      \[
        \begin{aligned}
          u = & \mathrm{e}^{-\int\left(-\frac{1}{x}\right)\mathrm{d}x}\left(C+\int\left(-\ln x\right)\mathrm{e}^{\int\left(-\frac{1}{x}\right)\mathrm{d}x}\mathrm{d}x\right) \\
          =   & x\left[C-\int\ln x\cdot\mathrm{d}\left(\ln x\right)\right] = x\left(C-\frac{1}{2}\ln^{2} x\right).
        \end{aligned}
      \]

      所以原方程通解为$xy\left(C-\frac{1}{2}\ln^{2} x\right)=1$.
  \end{enumerate}

  \section{综合题(每小题 7 分，共 14 分)}
  \begin{enumerate}
    \item[17.] \textbf{Solution}. 当$x\in(-1,0)\cup(0,+\infty)$时，$\ln f(x) = \frac{\ln(1+x)}{x}$，

      所以$f'(x)=f(x)\cdot \frac{\frac{x}{1+x}-\ln(1+x)}{x^{2}}=(1+x)^{\frac{1}{x}}
      \cdot\frac{x-(1+x)\ln(1+x)}{x^{2}(1+x)}$.

      当$x=0$时，
      \[
        \begin{aligned}
          f'(0)=\lim_{x\to0}\frac{f(x)-f(0)}{x}= & \lim_{x\to0}\frac{(1+x)^{\frac{1}{x}}-\mathrm{e}}{x}                                                                           \\
          =                                      & \,\mathrm{e}\cdot\lim_{x\to0}\frac{\mathrm{e}^{\frac{\ln(1+x)}{x}-1}-1}{x}= \mathrm{e}\cdot \lim_{x\to0}\frac{\ln(1+x)-x}{x^2} \\
          =                                      & \mathrm{e}\cdot \lim_{x\to0}\frac{x-\frac{1}{2}x^2+o(x^2)-x}{x^2}=-\frac{\mathrm{e}}{2}.
        \end{aligned}
      \]

      又
      \[
        \begin{aligned}
          \lim_{x\to0}f'(x) = & \lim_{x\to0}(1+x)^{\frac{1}{x}}\cdot\frac{x-(1+x)\ln(1+x)}{x^{2}(1+x)}            \\
          =                   & \,\mathrm{e}\cdot\lim_{x\to0}\frac{x-(1+x)\ln(1+x)}{x^{2}}                        \\
          =                   & \,\mathrm{e}\cdot\lim_{x\to0}\frac{1-\ln(1+x)-1}{2x}=-\frac{\mathrm{e}}{2}=f'(0),
        \end{aligned}
      \]

      所以$f'(x)$在点$x=0$处连续. 当$x\in(-1,0)\cup(0,+\infty)$时，$f'(x)$显然连续，

      故$f'(x)$在$(-1,+\infty)$上处处连续.

    \item[18.] \textbf{Solution}.旋转体的体积$V=$
      \[
        \begin{aligned}
          \int_{0}^{\frac{\pi}{2}}2\pi x f(x)\mathrm{d}x = & 2\pi\int_{0}^{\frac{\pi}{2}}x\cos x\mathrm{d}x                                                                                               \\
          =                                                & 2\pi\int_{0}^{\frac{\pi}{2}}x\mathrm{d}\sin x = 2\pi\left(x\sin x\bigg|_{0}^{\frac{\pi}{2}}-\int_{0}^{\frac{\pi}{2}}\sin x\mathrm{d}x\right) \\
          =                                                & 2\pi\left(\frac{\pi}{2}-1\right)=\pi^{2}-2\pi.
        \end{aligned}
      \]
  \end{enumerate}

  \section{证明题(每小题 5 分，共 10 分)}
  \begin{enumerate}
    \item[19.] \textbf{Proof}. 由Taylor公式，
      \[
        \begin{gathered}
          f(b)=f(a)+f'(a)(b-a)+\frac{f''(\xi_{1})}{2}(b-a)^2, \\
          f(a)=f(b)+f'(b)(a-b)+\frac{f''(\xi_{2})}{2}(a-b)^2,
        \end{gathered}
      \]

      其中$\xi_{1},\xi_{2}$介于$a$与$b$之间. 上述两式相减得
      \[
        \begin{aligned}
          f(b)-f(a) = 0 = & f(a)-f(b) + (b-a)(f'(a)+f'(b)) + \frac{1}{2}(b-a)^{2}(f''(\xi_{1})-f''(\xi_{2})) \\
          =               & (b-a)(f'(a)+f'(b)) + \frac{1}{2}(b-a)^{2}(f''(\xi_{1})-f''(\xi_{2})).
        \end{aligned}
      \]

      所以
      \[
        \begin{aligned}
          |f'(a)+f'(b)| = & \frac{1}{2}(b-a)|f''(\xi_{1})-f''(\xi_{2})|     \\
          \leq            & \frac{1}{2}(b-a)(|f''(\xi_{1})|+|f''(\xi_{2})|) \\
          =               & M(b-a).
        \end{aligned}
      \]

    \item[20.] \textbf{Solution}. 令$F(t)=\int_{a}^{t}f(x)\mathrm{d}x-\frac{(t-a)(f(a)+f(t))}{2}
      (x\geq a)$，

      则$F'(t)=f(t)-\frac{f(a)+f(t)+(t-a)f'(t)}{2}$，

      $F''(t)=f'(t)-\frac{1}{2}f'(t)-\frac{1}{2}\left[f'(t)+(t-a)f''(t)\right]=-\frac{1}{2}
      (t-a)f''(t)\geq 0$.

      所以$F'(t)$在$[a,+\infty)$上单调递增，$F'(t)\geq F'(a)=0$，

      所以$F(t)$在$[a,+\infty)$上单调递增，$F(t)\geq F(a)=0$，

      因此$\int_{a}^{b}f(x)\mathrm{d}x \geq \frac{(b-a)(f(a)+f(b))}{2}$.
  \end{enumerate}
\end{document}
